\subsection{Two auxiliary-field mechanisms in Yang--Mills theory}
\label{pg:sec:ym}

The finite-dimensional mirror of the Hubbard--Stratonovich step is the quartic Yang--Mills vertex, which can be made cubic with an auxiliary field in at least two ways.
The two mechanisms eliminate to the same quartic but differ in the anatomy of the auxiliary, and the comparison is what isolates the meaning of \enquote{invertible kinetic operator} for $p_{2,2}$.
We work both out, then compare.

\subsubsection{The color--kinematic mechanism}
\label{pg:sec:ym:ck}

In the color--kinematics-dual form of the Yang--Mills action the quartic vertex is written non-locally \cite{Borsten:2021:DoubleCopyHomotopy},
\begin{equation}
  \label{pg:eq:ym:quartic-nonlocal}
  \mathscr L^{(4)}
    \sim - \eta^{\mu\nu} \eta^{\kappa\rho} \eta^{\lambda\sigma} \,
      \pd^{12}_\mu \pd^{34}_\nu \,
      \frac{\tr\big( [A_\kappa, A_\lambda] [A_\rho, A_\sigma] \big)}{\square_{12}},
\end{equation}
where $\pd^{12}$ and $\square_{12}$ act on the first pair of fields.
The $1/\square$ is removed by a rank-three auxiliary field $G^a_{\mu\nu\kappa}$ with kinetic term $-\frac12 G\square G$ and a local coupling of the form $\pd A\,[A,A]$, giving a strictly cubic Lagrangian.
The mechanism is visible on a scalar representative with a single contraction, $G$ a scalar and $J$ standing for the local cubic current:
\begin{equation}
  \mathscr L(G)
    = - \frac12 \, G \square G
      - \frac{1}{\sqrt 2} \, G J .
\end{equation}
The equation of motion is $-\square G-J/\sqrt2=0$, so
\begin{equation}
  G
    = - \frac{1}{\sqrt 2} \, \frac{1}{\square} \, J,
\end{equation}
and substituting back, integrating by parts to let $\square$ act symmetrically,
\begin{equation}
  \mathscr L(G) \big|_{\rm on-shell}
    = - \frac12 \cdot \frac12 \, J \frac{1}{\square} J
      + \frac12 \, J \frac{1}{\square} J
    = \frac14 \, J \frac{1}{\square} J ,
\end{equation}
the non-local quartic with the correct coefficient.
The auxiliary kinetic operator is $\square$, invertible away from $p^2=0$.

\subsubsection{The ultralocal mechanism}
\label{pg:sec:ym:h}

The second mechanism introduces a Lie-algebra-valued antisymmetric auxiliary $h_{\mu\nu}=-h_{\nu\mu}$ with an algebraic kinetic term, no derivatives at all:
\begin{equation}
  \label{pg:eq:ym:Lh}
  \mathscr L_h
    = \frac14 \, \tr\big( h_{\mu\nu} h^{\mu\nu} \big)
      - \frac12 \, \tr\big( h_{\mu\nu} [A^\mu, A^\nu] \big).
\end{equation}
The $h$ equation of motion is algebraic,
\begin{equation}
  \frac{\pd \mathscr L_h}{\pd h_{\mu\nu}}
    = \frac12 \, h_{\mu\nu} - \frac12 \, [A_\mu, A_\nu]
    = 0
  \qquad\Longrightarrow\qquad
  \boxed{\; h_{\mu\nu} = [A_\mu, A_\nu] \;}
\end{equation}
so $h$ is exactly the commutator.
Substituting back,
\begin{equation}
  \begin{aligned}
    \mathscr L_h \big|_{h = [A, A]}
      &
      = \frac14 \, \tr\big( [A_\mu, A_\nu] [A^\mu, A^\nu] \big)
        - \frac12 \, \tr\big( [A_\mu, A_\nu] [A^\mu, A^\nu] \big)
    \\ &
      = - \frac14 \, \tr\big( [A_\mu, A_\nu] [A^\mu, A^\nu] \big),
  \end{aligned}
\end{equation}
the Yang--Mills quartic with the correct coefficient.
The kinetic operator is $\pd^2\mathscr L_h/\pd h^2=\frac12\,\mathbb 1$, a constant mass matrix, invertible everywhere with no momentum dependence.
Note the same $1:2$ magnitude ratio between the kinetic and the source term as in \eqref{p:eq:hs4:S4tilde}: at the saddle the two terms combine as $+1-2=-1$, in units of $\frac14\tr([A_\mu,A_\nu][A^\mu,A^\nu])$ --- the same doubling mechanism as the string action's $-1+2=+1$, with the overall sign flipped because the $h$-mechanism's kinetic term has the opposite sign from the $\Sigma$ one.
The elimination was checked symbolically with $\mathfrak{su}(2)$-valued $A_\mu$ in Euclidean signature; the identity is algebraic and signature-independent.

\subsubsection{The comparison}
\label{pg:sec:ym:compare}

\begin{center}
  \small
  \begin{tabular}{lccc}
    \toprule
    mechanism & auxiliary & kinetic operator & coupling \\
    \midrule
    color--kinematic & $G^a_{\mu\nu\kappa}$ (rank 3) & $\square$ (non-local) & $\pd A [A, A]$ \\
    ultralocal & $h_{\mu\nu}$ (rank 2) & $\frac12 \mathbb 1$ (ultralocal) & $h [A, A]$ \\
    closed string & $\Sigma \in \Sym^2 \mc H_0 / \ker \mc V_4$ & $\mc V_4$ (contact) & $\mc V_4(\Sigma, \Psi^2)$ \\
    \bottomrule
  \end{tabular}
\end{center}

What each mechanism buys is different.
The ultralocal one shows what \enquote{auxiliary} means at its simplest: an algebraic equation of motion, a constant invertible Hessian, no propagator, no question to ask.
The color--kinematic one shows that the kinetic operator of an auxiliary field can have a non-trivial inverse and that the substantive question is then its invertibility on the relevant space: for $G$, invertibility of $\square$ away from $p^2=0$.

The closed-string $\Sigma$ is like the $h$ form in that its kinetic operator $\mc V_4$ is of contact type, ultralocal in the moduli-space sense: the auxiliary propagates through a genuine contact vertex, with no $1/L_0^+$ and no propagator pole.
It is like the $G$ form in that invertibility is the substantive question: for $\Sigma$, invertibility of $\mc V_4$ on $\Sym^2\mc H_0$, or on the quotient by its kernel.
This is why the color--kinematic mechanism is the better model for the invertibility of $p_{2,2}$, and why it is kept rather than replaced by the simpler $h$ form: the $h$ form makes the trivially invertible end of the spectrum concrete, the $G$ form the non-trivially invertible end, and $\mc V_4$ sits between them, contact-type but with an invertibility that has to be established.

What does not carry over is the way the kernel is handled.
In Yang--Mills the operator $\square$ does have a kernel, the massless shell, and integrates out benignly through a Green's function because the degeneration is a measure-zero shell in a continuous momentum spectrum that diagonalizes $\square$.
The closed-string condition is an operator-kernel question on $\Sym^2\mc H_0$: there is no single invariant that diagonalizes the contact vertex on two-string states, and a normalizable symmetric null vector of $\mc V_4$ would be an honest flat direction of the action, not a removable shell.
The Yang--Mills analogy therefore motivates the expectation that $\mc V_4$ is non-degenerate on the symmetric square, and provides no argument for it; the quotient~\eqref{p:eq:hs4:quotient} is what makes the expectation unnecessary at this order.

